\section{Preliminaries}\label{prel}
Throughout $K$ will be a field. For a positive integer $n$ we denote by $U_n (K) = \{ \omega \in K \,|\, \omega^n = 1\}$ the cyclic group of $n$-th roots of unity in $K$, and by $\nu (n) = |U_n (K)|$ its order. Obviously $\nu (n)$ is a divisor of $n$; if $\nu (n) = n$, then any generator of $U_n (K)$ is called a~primitive $n$-th root of unity. In the sequel~$C_{n}$ will be the cyclic group of order $n$ generated by $g$ and $K[C_{n}]$ the corresponding group Hopf algebra. For any~$i,j \in \mathbb{N}$, $\delta_{i, j}$ denotes the Kronecker delta.

Let $m \in \mathbb{N}$, $m \geq 2$. Whenever we deal with Taft algebras of order $m^{2}$ we assume that the base field $K$ contains a~primitive $m$-th root of unity $q$. $T_{m^{2}}(q)$ denotes the Taft Hopf algebra of order $m$ over $K$, which is generated as an algebra by two elements~$h$ and~$x$ subject to the relations $h^{m} = 1$, $x^{m} = 0$ and $xh = qhx$. The coalgebra structure and antipode are given as follows
\begin{gather*}
\Delta(h) = h \otimes h, \qquad \Delta(x) = x \otimes h + 1 \otimes x, \qquad \epsilon(h) = 1,\\
 \epsilon(x) = 0,\qquad S(h) = h^{-1},\qquad S(x) = -xh^{-1}.
\end{gather*}
that is, $h$ is a group like element while $x$ is $(h,1)$-primitive. Sweedler's Hopf algebra is obtained by considering $m=2$ and $q = -1$. It can be easily checked that $\{h^{i}x^{j}\}_{0 \leq i, j \leq m-1}$ is a $K$-linear basis of the Taft algebra, the set of group-like elements is
$\mathcal{G}\big(T_{m^{2}}(q)\big) = \{1, h, h^{2},\dots, h^{m-1}\}$ and the primitive elements $\mathcal{P}_{h^{j}, 1}\big(T_{m^{2}}(q)\big)$ are given as follows for any $j = \{0, 1,\dots, m-1\}$
\begin{gather*}\mathcal{P}_{h^{j}, 1} \big(T_{m^{2}}(q)\big) = \begin{cases}
\alpha\big(h^{j} - 1\big), & \text{if } j \neq 1,\\
\beta(h-1)+ \gamma x, & \text {if } j=1,
\end{cases}\qquad \text{for some} \ \alpha, \beta, \gamma \in K.\end{gather*}
Unless specified otherwise, all algebras, coalgebras, bialgebras, Hopf algebras, tensor products and homomorphisms are over~$K$. For a coalgebra $(C, \Delta, \varepsilon)$, we use Sweedler's $\Sigma$-notation: $\Delta(c) = c_{(1)}\otimes c_{(2)}$, $(I\otimes \Delta)\Delta(c) = c_{(1)}\otimes
c_{(2)}\otimes c_{(3)}$, etc.\ (summation understood). Let $A$ and $H$ be two Hopf algebras. $H$ is called a right $A$-module coalgebra if $H$ is a coalgebra in the monoidal category ${\mathcal M}_A $ of right $A$-modules, i.e., there exists $\triangleleft \colon H \otimes A \rightarrow H$ a morphism of coalgebras such that $(H,\triangleleft) $ is a right $A$-module. A morphism between two right $A$-module coalgebras $(H, \triangleleft)$ and $(H',\triangleleft')$ is a morphism of coalgebras $\psi\colon H \to H'$ that is also right $A$-linear. Furthermore, $\psi$~is called unitary if
$\psi (1_H) = 1_{H'}$. Similarly, $A$ is a left $H$-module coalgebra if $A$ is a coalgebra in the monoidal category of left $H$-modules, that is there exists $\triangleright \colon H \otimes A \to A$ a morphism of coalgebras such that $(A, \triangleright)$ is also a~left $H$-module. The actions $\triangleleft \colon H \otimes A \rightarrow H$, $\triangleright: H \otimes A \rightarrow A$ are called trivial if $y \triangleleft a = \varepsilon_A (a) y$ and respectively $y \triangleright a = \varepsilon_H(y) a$, for all $a\in A$ and $y\in H$.

A Hopf algebra $E$ \emph{factorizes} through two Hopf algebras $A$ and $H$ if there exist injective Hopf algebra maps $i \colon A \to E $
and $j \colon H\to E$ such that the following map is bijective:
\begin{gather*}
A \otimes H \to E,\qquad a \otimes y \mapsto i(a) j(y).
\end{gather*}
A crucial result which characterizes Hopf algebras that factorize through two given Hopf algebras in terms of bicrossed products was proved by Majid (see for instance \cite[Theorem~7.2.3]{majid2}). We recall briefly the construction of the \emph{bicrossed product} of two Hopf algebras. It was introduced by Majid in \cite[Proposition 3.12]{majid} under the name of \emph{double cross product}. Throughout we shall adopt the name of bicrossed product from \cite[Theorem~IX 2.3]{Kassel} since it has also been used for similar constructions in many other fields (see, e.g., the group case \cite{Takeuchi}). The main ingredient for constructing bicrossed products are the so-called matched pairs. A \textit{matched pair} of Hopf algebras is a~quadruple $(A, H, \triangleleft, \triangleright)$, where $A$ and $H$ are Hopf algebras, $\triangleleft \colon H \otimes A \rightarrow H$,
$\triangleright\colon H \otimes A \rightarrow A$ are coalgebra maps such that $(A, \triangleright)$ is a left $H$-module coalgebra, $(H, \triangleleft)$ is a~right $A$-module coalgebra and the following compatibilities hold for any $a, b\in A$, $y, z\in H$
\begin{gather}
y \triangleright 1_{A} {=} \varepsilon_{H}(h)1_{A}, \qquad 1_{H} \triangleleft a = \varepsilon_{A}(a)1_{H}, \label{mp1} \\
y \triangleright(ab) {=} (y_{(1)} \triangleright a_{(1)}) \big( \big(y_{(2)}\triangleleft a_{(2)}\big)\triangleright b \big),\label{mp2} \\
(y z) \triangleleft a {=} \big( y \triangleleft \big(z_{(1)}\triangleright a_{(1)}\big) \big) (z_{(2)} \triangleleft a_{(2)}),\label{mp3} \\
y_{(1)} \triangleleft a_{(1)} \otimes y_{(2)} \triangleright a_{(2)} {=} y_{(2)} \triangleleft a_{(2)} \otimes y_{(1)}\triangleright a_{(1)}.\label{mp4}
\end{gather}
If $(A, H, \triangleleft, \triangleright)$ is a matched pair of Hopf algebras then the corresponding \textit{bicrossed product of Hopf algebras} $A \bowtie H$ is the tensor coalgebra $A\otimes H$ with the multiplication and antipode given as follows
\begin{gather*}
(a \bowtie y) \cdot (b \bowtie z) = a \big(y_{(1)}\triangleright b_{(1)}\big) \bowtie \big(y_{(2)} \triangleleft b_{(2)}\big) z, \\
S_{A \bowtie H} ( a \bowtie y ) = S_H \big(y_{(2)}\big) \triangleright S_A \big(a_{(2)}\big) \bowtie S_H \big(y_{(1)}\big) \triangleleft S_A \big(a_{(1)}\big)
\end{gather*}
for all $a,b \in A$, $y,z\in H$, where we denote $a\otimes y$ by $a\bowtie y$.

Although there are many interesting examples of bicrossed products such as the Drinfel'd double \cite[Theorem~IX.3.5]{Kassel} or the generalized quantum double \cite[Example~7.2.6]{majid} we will restrict to presenting only one example which is important for our purposes, namely the semi-direct (smash) product of Hopf algebras.

\begin{Example} \label{exempleban}
Let $(A, \triangleright)$ be a left $H$-module coalgebra and consider $H$ as a right $A$-module coalgebra via the trivial action, i.e., $y \triangleleft a = \varepsilon_A(a) y$. Then $(A,H, \triangleleft, \triangleright)$ is a matched pair of Hopf algebras if and only if $(A, \triangleright)$ is also a left $H$-module algebra and the following compatibility condition holds for all $y \in H$ and $a\in A$
\begin{gather*}
y_{(1)} \otimes y_{(2)} \triangleright a = y_{(2)} \otimes y_{(1)} \triangleright a.
\end{gather*}
In this case, the associated bicrossed product $A\bowtie H = A\# H$ is the left version of the \emph{semi-direct $($smash$)$ product} of Hopf algebras as defined by Molnar~\cite{Mo} in the cocommutative Hopf algebra setting where the compatibility condition~\eqref{mp4} holds automatically. Thus, $A\# H$ is the tensor coalgebra $A\otimes H$, with the following multiplication
\begin{gather*}
(a \# y) \cdot (b \# z):= a \big(y_{(1)} \triangleright b \big) \# y_{(2)} z
\end{gather*}
for all $a, b\in A$, $y, z\in H$, where we denote $a\otimes y$ by $a\# y$. Similarly one can define the right version of the smash product of Hopf algebras by considering $A$ as a left $H$-module coalgebra via the trivial action, i.e., $y \triangleright a := \varepsilon_H(y) a$.
\end{Example}

The next result, proved by Majid, establishes the connection between Hopf algebras which factorize through two given Hopf algebras and bicrossed products.

\begin{Theorem}[{\cite[Theorem 7.2.3]{majid2}}]\label{carMaj}
Consider two Hopf algebras $A$ and $H$. A Hopf algebra $E$ factorizes through $A$ and $H$ if and only if there exists a matched pair of Hopf algebras $(A, H, \triangleleft, \triangleright)$ and an isomorphism of Hopf algebras $E \cong A \bowtie H$.
\end{Theorem}
In light of Theorem~\ref{carMaj} the factorization problem for Hopf
algebras can be restated in terms of matched pairs and bicrossed products as follows:

\textit{For two given Hopf algebras $A$ and $H$, describe the set of all matched pairs $(A, H, \triangleright, \triangleleft)$ and classify
up to an isomorphism the corresponding bicrossed products $A \bowtie H$.}

The strategy we employ for classifying semi-direct products will rely on the following result which is a special case of \cite[Corollary~3.3]{abm}.

\begin{Theorem}\label{toatemorf}
Consider $A \# H$ and $A \#' H$ to be two semi-direct products associated to the left actions $\triangleright \colon H\otimes A \to A$ and
respectively $\triangleright' \colon H\otimes A \to A$. Then there exists a~bijective correspondence between the set of all morphisms of Hopf
algebras $\psi \colon A \# H \to A \# ' H $ and the set of all quadruples $(u, p, r, v)$, consisting of two unitary coalgebra maps $u\colon A \to A$ and $r\colon H \rightarrow A$, and two Hopf algebra maps $p\colon A \to H$ and $v\colon H \rightarrow H$ subject to the following compatibilities
\begin{gather}
u\big(a_{(1)}\big) \otimes p\big(a_{(2)}\big) = u\big(a_{(2)}\big) \otimes p\big(a_{(1)}\big),\label{C1}\\
r\big(t_{(1)}\big) \otimes v\big(t_{(2)}\big) = r\big(t_{(2)}\big) \otimes v\big(t_{(1)}\big),\label{C2}\\
u(ab) = u\big(a_{(1)}\big) \big( p \big(a_{(2)}\big) \triangleright' u(b) \big),\label{C3}\\
r(tw) = r\big(t_{(1)}\big) \big(v\big(t_{(2)}\big) \triangleright' r(w)\big),\label{C4}\\
r\big(t_{(1)}\big) \big(v\big(t_{(2)}\big) \triangleright' u(b) \big) = u \big(t_{(1)} \triangleright b_{(1)}\big) \big( p \big(t_{(2)}
\triangleright b_{(2)}\big) \triangleright' r\big(t_{(3)}\big) \big) ,\label{C5}\\
v(t) p (b)= p \big(t_{(1)} \triangleright b\big) v \big(t_{(2)}\big)\label{C6}
\end{gather}
for all $a, b \in A$, $t, w \in H$.

Under the above correspondence the morphism of Hopf algebras $\psi\colon A \# H \to A \# ' H $ corresponding to $(u, p, r, v)$ is given by
\begin{gather*}%\label{morfbicros}
\psi_{(u, p,r, v)}(a \# t) = u\big(a_{(1)}\big) \big( p\big(a_{(2)}\big) \triangleright' r\big(t_{(1)}\big) \big) \#' p\big(a_{(3)}\big) v\big(t_{(2)}\big)
\end{gather*}
for all $a \in A$ and $t\in H$.
\end{Theorem}

Finally, for the convenience of the reader we include here a straightforward result which will be intensively used in computing all matched pairs between $T_{m^{2}}(q)$ and $K[C_{n}]$.
\begin{Lemma}[{\cite[Lemma 4.1]{abm}}]\label{primitive}
Let $(A, H, \triangleleft, \triangleright)$ be a matched pair of Hopf algebras, $a, b \in \mathcal{G}(A)$ and $t, w\in \mathcal{G}(H)$. Then
\begin{enumerate}\itemsep=0pt
\item[$(1)$] $t \triangleright a \in \mathcal{G}(A)$ and $t\triangleleft a \in \mathcal{G}(H)$;

\item[$(2)$] if $c \in \mathcal{P}_{a, b}(A)$, then $t \triangleleft c \in \mathcal{P}_{t\triangleleft a, t\triangleleft b} (H)$ and
$t \triangleright c \in \mathcal{P}_{t\triangleright a, t\triangleright b} (A)$;

\item[$(3)$] if $z \in \mathcal{P}_{t, w}(H)$, then $z \triangleleft a \in \mathcal{P}_{t\triangleleft a, w\triangleleft a} (H)$ and $z \triangleright a \in \mathcal{P}_{t\triangleright a, w\triangleright a} (A)$.
\end{enumerate}
In particular, if $c$ is an $(a, 1)$-primitive element of $A$, then $t \triangleright c$ is an $(t\triangleright a, 1)$-primitive element of $A$ and $t \triangleleft c$ is an $(t\triangleleft a, 1)$-primitive element of~$H$.
\end{Lemma}

\section{Main results}\label{2}

In this section we deal with the description and classification of Hopf algebras which factorize through $T_{m^{2}}(q)$ and $K[C_{n}]$ for any $m, n \in \mathbb{N}$, $n \geqslant 2$, $m \geqslant 3$. The case $m = 2$ was recently considered in~\cite{abm}. Throughout, the field $K$ will be assumed to contain a primitive $m$-th root of unity $q$. We start by describing all possible matched pairs $(T_{m^{2}}(q), K[C_n], \triangleleft,
\triangleright)$: exactly as in~\cite{abm}, it turns out that the aforementioned matched pairs are in bijection with the cyclic group~$U_n (K)$ of
$n$-th roots of unity in~$K$. More precisely, the right action $\triangleleft \colon K[C_{n}] \otimes T_{m^{2}}(q) \to K[C_{n}]$ is trivial while
the left action $\triangleright \colon K[C_{n}] \otimes T_{m^{2}}(q) \to T_{m^{2}}(q)$ is implemented by an $n$-th root of unity as described below.

\begin{Proposition} \label{mapex1}
There exists a bijective correspondence between the set of all matched pairs $(T_{m^{2}}(q), K[C_n], \triangleleft, \triangleright)$ and $U_n (K)$. More precisely, the matched pair $(\triangleleft, \triangleright)$ corresponding to an $n$-th root of unity $\omega \in U_n (K)$ is given by
\begin{gather}\label{mp}
g^{i} \triangleright h^{j} x^{k} = \omega^{ik} h^{j} x^{k}, \qquad g^{i} \triangleleft h^{j} x^{k} = g^{i} \varepsilon\big(x^{k}\big)
\end{gather}
for all $i = 0,\dots, n-1$ and $j, k = 0,\dots, m-1$.
\end{Proposition}

\begin{proof} Let $(T_{m^{2}}(q), K[C_n], \triangleleft, \triangleright)$ be a~matched pair. We start by proving that $g \triangleright h = h$.
Indeed, by Lemma~\ref{primitive} we have $g \triangleright h \in \mathcal{G}\big(T_{m^{2}}(q)\big) = \{1, h,\dots,h^{m-1}\}$. Suppose first that $g \triangleright h = 1$. Then, by induction we have $1 = g^{n} \triangleright h = h$ which is obviously a contradiction. Assume now that $g \triangleright h = h^{t}$, where $t \in \{2, 3,\dots, m-1\}$. Now Lemma~\ref{primitive} implies $g \triangleright x \in \mathcal{P}_{h^{t}, 1}\big(T_{m^{2}}(q)\big)$. Then, as $t \neq 1$ we obtain $g \triangleright x = \alpha(1-h^{t})$ for some $\alpha \in K$. Hence, $g^{2} \triangleright x = g \triangleright (g \triangleright x) = \alpha - \alpha g \triangleright h^{t}$. Again by induction we arrive at $x = g^{n} \triangleright x = \alpha - \alpha g^{n-1} \triangleright h^{t}$ which leads to another contradiction since Lemma~\ref{primitive} yields $g^{n-1} \triangleright h^{t} \in \mathcal{G}\big(T_{m^{2}}(q)\big) = \{1, h,\dots,h^{m-1}\}$. Hence $g \triangleright h = h$ as desired. Furthermore, this
also implies that $g \triangleright x = \alpha(1-h)+ \beta x$ for some~$\alpha, \beta \in K$. Using induction again we obtain
\begin{gather*}
g^{n} \triangleright x = \alpha\big(1 + \beta + \cdots + \beta^{n-1}\big) (1-h) + \beta^{n} x.
\end{gather*}
As we also have $g^{n} \triangleright x = x$ it follows that
\begin{gather*}
\beta^{n} = 1, \qquad \alpha\big(1 + \beta + \cdots + \beta^{n-1}\big)=0.
\end{gather*}
Next we look at the right action $\triangleleft$. Again by Lemma~\ref{primitive} we have $g \triangleleft h \in \mathcal{G}\big(K[C_n]\big) = \{1, g,\dots, g^{n-1}\}$, therefore $g \triangleleft h = g^{t}$ for some $t \in \{0,1, \dots, n-1\}$. If $g \triangleleft h = 1$ we obtain by induction
$1 = g \triangleleft h^{n} = g \triangleleft 1 = g$ which is a~contradiction. Thus $g \triangleleft h = g^{t}$ for some $t \in \{1, 2,\dots, n-1\}$. Moreover, we also have $g \triangleleft x \in \mathcal{P}_{g^{t}, g} \big(K[C_{n}]\big)$, i.e., $g \triangleleft x = \mu(g - g^{t})$ for some $\mu \in K$. Applying the compatibility condition~\eqref{mp4} for the pair $(g, x)$ yields
\begin{gather*}
\mu g \otimes h - \mu g^{t} \otimes h + \alpha g \otimes 1 - \alpha g \otimes h + \beta g \otimes x \\
\qquad{} = \alpha g^{t} \otimes 1 - \alpha g^{t} \otimes h + \beta g^{t} \otimes x + \mu g \otimes 1 - \mu g^{t} \otimes 1.
\end{gather*}
Now if $t \neq 1$ then we must have $\beta = 0$ which again contradicts $\beta^{n} = 1$. Hence $t = 1$ which implies $g \triangleleft h = g$ and respectively $g \triangleleft x = 0$. Moreover, the compatibility condition~\eqref{mp4} is now trivially fulfilled.

Next we make use of the compatibility condition~\eqref{mp2}. More precisely we have
\begin{gather*}
g \triangleright hx \stackrel{\eqref{mp2}} {=} (g \triangleright h) \big((g \triangleleft h) \triangleright x\big)
= h (g \triangleright x) = \alpha h - \alpha h^{2} + \beta hx,\\
g \triangleright xh \stackrel{\eqref{mp2}} {=} \big(g \triangleright x_{(1)}\big) \big(\big(g \triangleleft x_{(2)}\big) \triangleright h\big)
= (g \triangleright x) \big((g \triangleleft h) \triangleright h\big) = \alpha h - \alpha h^{2} + \beta xh.
\end{gather*}
As $xh = q hx$, putting all the above together gives
\begin{gather*}
q\alpha h - q\alpha h^{2} + q\beta hx = \alpha h - \alpha h^{2} + \beta xh,
\end{gather*}
and since $q \neq 1$ we obtain $\alpha = 0$. Thus $g \triangleright x = \beta x$. To summarize, for all $i \in \{0, 1,\dots, n-1\}$ and $j \in \{0, 1,\dots, m-1\}$ we have
\begin{gather*}
g \triangleleft h^{j} = g, \qquad g \triangleleft x^{j} = g \varepsilon\big(x^{j}\big), \qquad g^{i} \triangleright h = h, \qquad g^{i} \triangleright x = \beta^{i} x,\qquad {\rm where} \quad \beta \in U_{n}(K).
\end{gather*}
Now using induction and the compatibility conditions~\eqref{mp2} and respectively~\eqref{mp3} we obtain the formulae in~\eqref{mp}. To start with, using~\eqref{mp3} we obtain $g^{i} \triangleleft h = g^{i}$ for all $i \in \{0, 1,\dots, n-1\}$. Indeed, we have
\begin{gather*}
g^{2} \triangleleft h \stackrel{\eqref{mp3}} = \big(g \triangleleft (g \triangleright h)\big) (g \triangleleft h)= g^{2},\\
g^{3} \triangleleft h \stackrel{\eqref{mp3}} = \big(g^{2} \triangleleft (g \triangleright h)\big) (g \triangleleft h)= \big(g^{2} \triangleleft h\big)(g \triangleleft h) = g^{3},
\end{gather*}
and by induction we arrive at the desired conclusion. Furthermore, as $\triangleleft$ is a right action of $T_{m^{2}}(q)$ on $K[C_{n}]$, we obtain
\begin{gather}\label{001}
g^{i} \triangleleft h^{j} = g^{i},\qquad \text{for all} \quad i \in \{0, 1,\dots, n-1\} \quad {\rm and}\quad j \in \{0, 1,\dots, m-1\}.
\end{gather}
Now induction together with~\eqref{mp2} and~\eqref{001} yields
\begin{gather}\label{002}
g^{i} \triangleright h^{j} = h^{j}, \qquad \text{for all} \quad i \in \{0, 1,\dots, n-1\} \quad {\rm and} \quad j \in \{0, 1,\dots, m-1\}.
\end{gather}
Next we use~\eqref{mp3} in order to prove that $g^{i} \triangleleft x =0$. Indeed, we have
\begin{gather*}
g^{2} \triangleleft x \stackrel{\eqref{mp3}} {=} \big(g \triangleleft \big(g \triangleright x_{(1)}\big)\big) \big(g \triangleleft
x_{(2)}\big) = \big(g \triangleleft (g \triangleright x)\big) (g \triangleleft h) + \big(g \triangleleft (g \triangleright 1)\big) (g \triangleleft x)=0, \\
g^{3} \triangleleft x \stackrel{\eqref{mp3}} {=} \big(g^{2} \triangleleft \big(g \triangleright x_{(1)}\big)\big) \big(g \triangleleft
x_{(2)}\big)= \big(g^{2} \triangleleft (g \triangleright x)\big) (g \triangleleft h) + \big(g^{2} \triangleleft (g \triangleright 1)\big) (g \triangleleft x)=0,
\end{gather*}
and the conclusion follows again by induction. Moreover, we actually have
\begin{gather*}%\label{003}
g^{i} \triangleleft x^{j} = g^{i} \varepsilon\big(x^{j}\big), \qquad \text{for all} \quad i \in \{0, 1,\dots, n-1\} \quad {\rm and} \quad j \in \{0, 1,\dots, m-1\}.
\end{gather*}
Finally, we will use~\eqref{mp3} in order to prove that
\begin{gather}\label{003}
g^{i} \triangleright x^{j} = \beta^{ij}x^{j}, \qquad \text{for all} \quad i \in \{0, 1,\dots, n-1\} \quad {\rm and} \quad j \in \{0, 1,\dots, m-1\}.
\end{gather}
To this end, we have
\begin{gather*}
g^{i} \triangleright x^{2} \stackrel{\eqref{mp2}} {=} \big(g^{i} \triangleright x_{(1)}\big)\big( \big(g^{i} \triangleleft
x_{(2)}\big)\triangleright x\big) \\
\hphantom{g^{i} \triangleright x^{2}}{} \ = \big(g^{i} \triangleright x\big)\big( \big(g^{i} \triangleleft h\big)\triangleright x\big) + \big(g^{i} \triangleright 1\big)\big( \big(g^{i} \triangleleft x\big)\triangleright x\big) = \beta^{2i}x^{2},\\
g^{i} \triangleright x^{3} \stackrel{\eqref{mp2}} {=} \big(g^{i} \triangleright x_{(1)}\big)\big( \big(g^{i} \triangleleft
x_{(2)}\big)\triangleright x^{2}\big)\\
 \hphantom{g^{i} \triangleright x^{3}}{} \ = \big(g^{i} \triangleright x\big)\big( \big(g^{i} \triangleleft h\big)\triangleright x^{2}\big) + \big(g^{i} \triangleright 1\big)\big( \big(g^{i} \triangleleft x\big)\triangleright x^{2}\big) =\beta^{3i}x^{3},
\end{gather*}
and as before the conclusion follows by induction. Putting all the above together we obtain
\begin{gather*}
g^{i} \triangleright h^{j} x^{k} \stackrel{\eqref{mp2}} {=} \big(g^{i} \triangleright h^{j}\big)\big( \big(g^{i} \triangleleft h^{j}\big)\triangleright x^{k}\big) \stackrel{\eqref{001}, \, \eqref{002}} {=} h^{j} \big(g^{i} \triangleright x^{k}\big)\stackrel{\eqref{003}} {=} \beta^{ik} h^{j} x^{k}.
\end{gather*}
Similarly one can prove the second part of \eqref{mp} and the proof is now finished.
\end{proof}

Next we show that a Hopf algebra $E$ factorizes through $T_{m^{2}}(q)$ and $K[C_n]$ if and only if $E \cong T_{nm^{2}}^{\omega}(q)$, where $T_{nm^{2}}^ {\omega}(q)$ is a quantum group associated to an $n$-th root of unity $\omega$ as depicted below.

\begin{Corollary}\label{primaclasi}
A Hopf algebra $E$ factorizes through $T_{m^{2}}(q)$ and $K[C_n]$ if and only if $E \cong T_{nm^{2}}^ {\omega}(q)$, for some $\omega \in U_n (K)$, where we denote by $T_{nm^{2}}^{\omega}(q)$ the Hopf algebra generated by $g$, $h$ and $x$ subject to the relations
\begin{gather*}
g^{n} = h^{m} = 1,\qquad x^{m} = 0,\qquad x h = q h x, \qquad h g = g h,\qquad g x = \omega x g
\end{gather*}
with the coalgebra structure and antipode given by
\begin{gather*}
\Delta(g) = g \otimes g, \qquad \Delta(h) = h \otimes h, \qquad \Delta(x) = x \otimes h + 1 \otimes x,\\
\varepsilon(h) = \varepsilon(g) = 1, \qquad \varepsilon(x) = 0, \qquad S(h) = h^{m-1}, \qquad S(x) = -xh^{m-1}, \qquad S(g) = g^{n-1}.
\end{gather*}
\end{Corollary}

\begin{proof} By Theorem~\ref{carMaj} any Hopf algebra~$E$ which factorizes through $T_{m^{2}}(q)$ and $K[C_n]$ is isomorphic to a bicrossed product
between the aforementioned Hopf algebras. Furthermore, all bicrossed products between $T_{m^{2}}(q)$ and $K[C_n]$ are described in Proposition~\ref{mapex1}: the left action is completely determined by an $n$-th root of unity $\omega$ as in~\eqref{mp} while the right action is trivial. Therefore, $E$ is in fact isomorphic to a smash product $T_{m^{2}}(q) \# K[C_n]$ associated to the left action defined in~\eqref{mp}. Now, up to canonical identification, $T_{m^{2}}(q) \# K[C_n]$ is generated as an algebra by $h = h \# 1$, $x = x \# 1$ and $g = 1\# g$. Hence, we have
\begin{gather*}
gh = (1\# g) (h \# 1)= g \triangleright h \# g \triangleleft h = h \# g = (h \# 1)(1 \# g) = hg,\\
g x = (1 \# g) (x \# 1) = g \triangleright x_{(1)}
\# g \triangleleft x_{(2)} = g \triangleright x \# g \triangleleft h + g \triangleright 1 \# g \triangleleft x =\omega x \# g = \omega x g. \tag*{\qed}
\end{gather*}\renewcommand{\qed}{}
\end{proof}
\begin{Remark}Let $\omega \in U_{n}(K)$. It can be easily seen that a $K$-basis in $T_{nm^{2}}^ {\omega}$ is given by $\{ g^{i}, h^{j} g^{i},
 x^{k} g^{i}, h^{j} x^{k} g^{i} \,| \,i = 0, 1,\dots, n-1, \,{\rm and}\, j, k = 1, 2,\dots, m-1 \}$. Obviously $T_{nm^{2}}^ {\omega}$ is a Hopf algebra of dimension $nm^{2}$ over~$K$.
\end{Remark}

Our next result provides the classification of Hopf algebras which factorize through the Taft algebra and the group Hopf algebra $k[C_{n}]$. To start with, if the Hopf algebras $T_{nm^{2}}^{\omega}$ and $T_{kl^{2}}^ {\omega '}$ are isomorphic it is straightforward to see, by comparing the order of their group of group-like elements and dimensions, that $n=k$ and $m=l$. Therefore, we will focus on finding necessary and sufficient conditions for two Hopf algebras $T_{nm^{2}}^ {\omega}$ and $T_{nm^{2}}^ {\omega '}$ to be isomorphic.

\begin{Theorem}\label{izo} Let $\xi$ be a generator of $U_n (K)$ and $t, t' \in \{ 0, 1,\dots, \nu (n) -1 \}$. Then the Hopf algebras $T_{nm^{2}}^ {\xi^{t}}$ and $T_{nm^{2}}^ {\xi^{t'}}$ are isomorphic if and only if there exist $l \in \{ 0, 1,\dots, m-1 \}$, $s \in \{ 0, 1,\dots, n-1 \}$ such that $(s, n) = 1$ and $\xi^{st' - t} = q^{l}$.
\end{Theorem}

\begin{proof} We know from Corollary~\ref{primaclasi} that $T_{nm^{2}}^ {\xi^{t}}$ (respectively $T_{nm^{2}}^ {\xi^{t'}}$) is the smash product corresponding to the root of unity $\xi^{t}$ (respectively $\xi^{t'}$) as in Proposition~\ref{mapex1}. By Theorem~\ref{toatemorf} the set of all Hopf algebra morphisms between $T_{nm^{2}}^ {\xi^{t}}$ and $T_{nm^{2}}^ {\xi^{t'}}$ is parameterized by the quadruples $(u, p, r, v)$, where $u\colon T_{m^{2}}(q) \to T_{m^{2}}(q)$, $r\colon K[C_{n}] \to T_{m^{2}}(q)$ are unitary coalgebra maps and $p\colon T_{m^{2}}(q) \to K[C_{n}]$, $v\colon K[C_{n}] \to K[C_{n}]$ are Hopf algebra maps satisfying the compatibility conditions~\eqref{C1}--\eqref{C6}. Our aim is to describe completely the above quadruples for which the corresponding morphism of Hopf algebras $\psi_{(u,p,r,v)}$ is bijective. We start by describing the two Hopf algebra maps $p$ and $v$. Obviously, any Hopf algebra map $v\colon K[C_{n}] \to K[C_{n}]$ is completely determined by an integer $s \in \{ 0, 1, \dots, n-1 \}$ such that $v(g) = g^{s}$. Similarly, using Lemma~\ref{primitive} we get $p(h) = g^{c}$ for some $c \in \{ 0, 1, \dots, n-1 \}$ such that $n \,|\, mc$. Furthermore, $p(x) = \lambda(g^{c} - 1)$ for some scalar $\lambda \in K$. Now imposing the compatibility condition $p(xh) = q p(hx)$ to hold true we obtain $\lambda g^{c}(g^{c} - 1) (1-q) = 0$ and since $q \neq 1$ we are lead to $\lambda(g^{c} - 1) = 0$, that is $p(x) = 0$. Now we turn to the two coalgebra maps $u\colon T_{m^{2}}(q) \to T_{m^{2}}(q)$, $r\colon K[C_{n}] \to T_{m^{2}}(q)$. We obviously have $r(1) = 1$, $r(g) = h^{l}$ for some $l \in \{ 0, \dots, m-1 \}$. The compatibility condition~\eqref{C4} gives
\begin{gather*}
r\big(g^{2}\big) = r(g)\big(v(g) \triangleright ' r(g)\big) = h^{l} \big(g^{s} \triangleright ' h^{l}\big) = h^{2l}.
\end{gather*}
Using induction and~\eqref{C4} we get $r(g^{i}) = h^{il}$ for any $i \in \mathbb{N}$. In particular, we have $1 = r(g^{n}) = h^{nl}$ and therefore $m \,|\, nl$. Finally we are left to describe the coalgebra maps $u\colon T_{m^{2}}(q) \to T_{m^{2}}(q)$. We have $u(1) = 1$ and $u(h)=h^{d}$ for some $d \in \{ 0, 1,\dots, m-1 \}$. Now we make use of the compatibility condition~\eqref{C3}
\begin{gather*}
u\big(h^{2}\big) = u(h) \big(p(h) \triangleright ' u(h)\big) = h^{d} \big(g^{c} \triangleright ' h^{d}\big) = h^{2d}.
\end{gather*}
As before, using induction and~\eqref{C3} we get $u(h^{j}) = h^{jd}$ for any $j \in \mathbb{N}$. Lemma~\ref{primitive} gives $u(x) \in
\mathcal{P}_{h^{d}, 1} \big(T_{m^{2}}(q)\big) = \begin{cases}
\alpha(h^{d}-1), & \text{if } d \neq 1,\\
\alpha(h-1)+ \gamma x, & \text{if } d=1,
\end{cases}$ for some $\alpha, \gamma \in K$. Suppose first that $d \neq 1$ and thus $u(x) = \alpha(h^{d}-1)$. Then we have
\begin{gather*}
u(hx) \stackrel{\eqref{C3}} {=} u(h) \big(p(h) \triangleright ' u(x)\big) = h^{d} \big(g^{c} \triangleright ' \alpha\big(h^{d}-1\big)\big) = \alpha h^{d} \big(h^{d}-1\big),\\
u(xh) \stackrel{\eqref{C3}} {=} u\big(x_{(1)}\big) \big(p\big(x_{(2)}\big) \triangleright ' u(h)\big) = u(x) \big(p(h) \triangleright ' u(h)\big) + p(x)
\triangleright ' u(h) = \alpha h^{d} \big(h^{d}-1\big).
\end{gather*}
Now $u(xh) = q u(hx)$ gives $\alpha (q-1) h^{d}(h^{d}-1) = 0$ and since $q \neq 1$ we obtain $\alpha(h^{d}-1) = 0$. Therefore $u(x) = 0$. However, we will see that in this case the corresponding Hopf algebra morphism $\psi_{(u, p, r, v)}$ is not an isomorphism. Indeed we have
\begin{gather*}
\psi_{(u, p, r, v)}(x \# 1) = u\big(x_{(1)}\big) \big(p\big(x_{(2)}\big) \triangleright ' r(1)\big) \# p(x_{(3)}) v(1),\\
\hphantom{\psi_{(u, p, r, v)}(x \# 1)}{} = u\big(x_{(1)}\big) \# p\big(x_{(2)}\big)= u(x) \# p(h) + 1 \# p(x) = 0,
\end{gather*}
which is a contradiction. Therefore we must have $d = 1$ and $u(x) = \alpha (h-1) + \gamma x$. Hence, we obtain
\begin{gather*}
u(hx) \stackrel{\eqref{C3}} {=} u(h) \big(p(h) \triangleright '
u(x)\big) = h \big[g^{c} \triangleright ' \big(\alpha(h-1)+
\gamma x \big)\big] = \alpha h(h-1)+ \gamma \xi^{t'c} hx,\\
u(xh) \stackrel{\eqref{C3}} {=} u\big(x_{(1)}\big) \big(p\big(x_{(2)}\big)
\triangleright ' u(h)\big) = u(x) \big(p(h) \triangleright ' u(h)\big) + p(x) \triangleright ' u(h) \\
\hphantom{u(xh)}{} \ = \big[\alpha(h-1) + \gamma x\big](g^{c} \triangleright ' h) = \alpha h(h-1) + \gamma q hx.
\end{gather*}
Now $u(xh) = q u(hx)$ gives $\alpha = q \alpha$ and $q \gamma \xi^{t'c} = q \gamma$. As $q \neq 1$ we obtain $\alpha = 0$. Thus $\gamma \neq 0$ (otherwise we would have $u(x) = 0$ which leads to the same contradiction as in the previous case) and so $\xi^{t'c} = 1$. Hence, $u$ is given as follows
\begin{gather*}
u(h^{i}) = h^{i}, \qquad u\big(x^{i}\big) = \gamma^{i} x^{i}, \qquad \gamma \in K^{*}, \qquad i \in \{0, 1,\dots, m-1\}.
\end{gather*}
Putting all together, the quadruple $(u, p, r, v)$ is given as follows for all $i, j \in \mathbb{N}$
\begin{gather}
v\big(g^{i}\big) = g^{is}, \qquad {\rm with} \quad s \in \{ 0,\dots, n-1 \}, \nonumber\\ %\label{cuad1}\\
p\big(h^{i}x^{j}\big) = g^{ic} \delta_{0, j}, \qquad {\rm with} \quad c \in \{ 0,\dots, n-1 \} \quad \text{such that} \quad n \,|\, mc \quad {\rm and} \quad \xi^{t'c} = 1,\nonumber\\ %\label{cuad2}\\
r\big(g^{i}\big) = h^{il},\qquad {\rm with} \quad l \in \{ 0,\dots, m-1 \} \quad \text{such that} \quad m \,|\, nl,\label{cuad3}\\
u\big(h^{i}\big) = h^{i}, \qquad u\big(x^{i}\big) = \gamma^{i} x^{i}, \qquad {\rm with} \quad \gamma \in K^{*}. \nonumber %\label{cuad4}
\end{gather}
We are still left to check the compatibilities~\eqref{C1},~\eqref{C5} and~\eqref{C6}. Remark that~\eqref{C2} is trivially fulfilled as $K[C_{n}]$ is cocommutative.

\eqref{C1} is trivially fulfilled for $a = h$ while for $a = x$ comes down to the following
\begin{gather*}
 u\big(x_{(1)}\big) \otimes p\big(x_{(2)}\big) = u\big(x_{(2)}\big) \otimes p\big(x_{(1)}\big) \\
\qquad{} \Leftrightarrow \quad u(x) \otimes p(h) + u(1) \otimes p(x) = u(h) \otimes p(x) + u(x) \otimes p(1) \\
\qquad{} \Leftrightarrow \quad \gamma x \otimes g^{c} = \gamma x \otimes 1,
\end{gather*}
and since $\gamma \neq 0$ we obtain $c = 0$ and thus $p$ is the trivial map, i.e., $p(a) = \varepsilon(a) 1$ for all $a \in T_{m^{2}}(q)$. Now~\eqref{C6} is trivially fulfilled while~\eqref{C5} comes down to the following
\begin{gather*}
r\big(t_{(1)}\big) \big(v\big(t_{(2)}\big) \triangleright' u(b) \big) {=} u \big(t_{(1)} \triangleright b\big) r\big(t_{(2)}\big).
\end{gather*}
By setting $t = g$ and $b = x$ we obtain
\begin{gather*}
 r(g)\big(v(g) \triangleright' u(x) \big) = u(g \triangleright x) r(g)\quad
\!\Leftrightarrow\! \quad h^{l} (g^{s} \triangleright' \gamma x) = u\big(\xi^{t} x\big)h^{l}\quad
\!\Leftrightarrow\!\quad \gamma \xi^{st'} h^{l} x = \gamma \xi^{t} q^{l} h^{l} x,\!
\end{gather*}
and thus $\xi^{st' - t} = q^{l}$. This last equality together with $\nu(n) \,|\, n$ gives $m \,|\, nl$. Thus, the compatibility condition in~\eqref{cuad3} is trivially fulfilled.

To summarize, we have proven that if the Hopf algebras $T_{nm^{2}}^ {\xi^{t}}$ and $T_{nm^{2}}^ {\xi^{t'}}$ are isomorphic then there exist $l \in \{ 0, 1,\dots, m-1 \}$ and $s \in \{ 0, 1,\dots, n-1 \}$ such that $\xi^{st' - t} = q^{l}$. In order to complete the first part of the proof we are left to show that $(s, n) = 1$. To this end, $v$ is bijective as a consequence of $\psi$ being bijective. Indeed, if $\psi^{-1}$ is the inverse of $\psi$ then using the same arguments as above one can easily prove that there exists a unitary coalgebra map $\overline{r}\colon K[C_{n}] \to T_{m^{2}}(q)$ and a Hopf algebra map $\overline{v}\colon K[C_{n}] \to K[C_{n}]$ such that for any $z \in K[C_{n}]$ we have $\psi^{-1} (1 \# z) = \overline{r}(z_{(1)}) \# \overline{v}(z_{(2)})$. This gives
\begin{gather*}
1 \# z = \psi \circ \psi^{-1} (1 \# z) = u \big(\overline{r}\big(z_{(1)}\big)\big) r \big(\overline{v}\big(z_{(2)}\big)\big) \# v\big(\overline{v}\big(z_{(3)}\big)\big).
\end{gather*}
By applying $\varepsilon \otimes {\rm Id}_{K[C_{n}]}$ to the above identity yields $v \circ \overline{v}= {\rm Id}_{K[C_{n}]}$. Similarly one can prove that $\overline{v} \circ v= {\rm Id}_{K[C_{n}]}$ and we can conclude that $v$ is indeed an isomorphism. Obviously $v$ is bijective if and only if $(s, n) = 1$.

Assume now that there exist $l \in \{ 0, 1,\dots, m-1 \}$, $s \in \{ 0, 1,\dots, n-1 \}$ such that $\xi^{st' - t} = q^{l}$ and $(s, n) = 1$. Consider two unitary coalgebra maps $u_{\gamma}\colon T_{m^{2}}(q) \to T_{m^{2}}(q)$, $r_{l}\colon K[C_{n}] \to T_{m^{2}}(q)$ and a~Hopf algebra map $v_{s}\colon K[C_{n}] \to K[C_{n}]$ defined as follows for all $i, j \in \{0, 1,\dots, m-1\}$ and $k \in \{0, 1,\dots, n-1\}$
\begin{gather*}
u_{\gamma}\big(h^{i}x^{j}\big) = \gamma^{j} h^{i}x^{j}, \qquad {\rm where}\quad \gamma \in K^{*}, \qquad r_{l}\big(g^{k}\big) = h^{kl}, \qquad v_{s}\big(g^{k}\big) = g^{ks}.
\end{gather*}
We claim that the following map is an isomorphism of Hopf algebras
\begin{gather*}
\psi_{(u_{\gamma}, \varepsilon, r_{l}, v_{s})}\colon \ T_{nm^{2}}^ {\xi^{t}} \to T_{nm^{2}}^ {\xi^{t'}}, \qquad \psi_{(u_{\gamma}, \varepsilon, r_{l}, v_{s})}(a \# y) = u_{\gamma}(a) r_{l}\big(y_{(1)}\big) \# v_{s}\big(y_{(2)}\big).
\end{gather*}
To start with, $\psi_{(u_{\gamma}, \varepsilon, r_{l}, v_{s})}$ is a Hopf algebra map according to the first part of the proof. We only need to show that $\psi_{(u_{\gamma}, \varepsilon, r_{l}, v_{s})}$ is bijective as well. To this end, since $(s, n) = 1$, there exist $\tau,\mu \in \mathbb{Z}$ such that $s \tau + n \mu = 1$. Furthermore, there exist unique integers $\alpha$, $\beta$, $\tau_{1}$, $\tau_{2}$ such that
\begin{gather}\label{impr}
\tau = \alpha n + \tau_{1} \qquad {\rm and} \qquad l \tau = \beta m + \tau_{2},
\end{gather}
where $\tau_{1} \in \{0, 1,\dots, n-1\}$, $\tau_{2} \in \{0, 1,\dots, m-1\}$. It will be worth our while to point out the following easy consequence of~\eqref{impr}
\begin{gather}\label{impr2}
l\tau_{1} = \beta m - \alpha nl + \tau_{2}.
\end{gather}
As $(\tau, n) = 1$ we obviously have $(\tau_{1}, n) = 1$ as well. We will prove that we also have $\xi^{\tau_{1} t - t'} = q^{m-\tau_{2}}$. First recall that, as noticed before, $\xi^{st' - t} = q^{l}$ implies $m \,|\, nl$. Now employing the equalities $\xi^{st' - t} = q^{l}$ and $s \tau + n \mu = 1$ we obtain
\begin{gather*}
\xi^{\tau_{1} t - t'} = \big(\xi^{t}\big)^{\tau_{1}} \xi ^{-t'} = \big(\xi^{st'} q^{-l}\big)^{\tau_{1}} \xi ^{-t'} =
\xi^{(s\tau_{1} -1)t'} q^{-l\tau_{1}} \stackrel{\eqref{impr}} {=} \xi^{(s\tau - s \alpha n -1)t'} q^{-l\tau_{1}} \\
\hphantom{\xi^{\tau_{1} t - t'}}{} = \xi^{(s\tau -1)t'} q^{-l\tau_{1}}= \xi^{-n \mu t'} q^{-l\tau_{1}} = q^{-l\tau_{1}} \stackrel{\eqref{impr2}} {=} q^{-\beta m + \alpha nl - \tau_{2}} \stackrel{m \,|\, nl} {=} q^{- \tau_{2}} = q^{m - \tau_{2}}.
\end{gather*}
Now since $\gamma^{-1} \in k^{*}$, $\tau_{1} \in \{0, 1,\dots, n-1\}$, $m - \tau_{2} \in \{0, 1,\dots, m-1\}$ and $\xi^{\tau_{1} t - t'} = q^{m-\tau_{2}}$, it follows from the first part of the proof that the map $\psi_{(u_{\gamma^{-1}}, \varepsilon, r_{m-\tau_{2}}, v_{\tau_{1}})}\colon T_{nm^{2}}^ {\xi^{t'}} \to T_{nm^{2}}^ {\xi^{t}}$ defined below is a Hopf algebra morphism
\begin{gather*}
\psi_{(u_{\gamma^{-1}}, \varepsilon, r_{m-\tau_{2}}, v_{\tau_{1}})}(a \# y) = u_{\gamma^{-1}}(a) r_{m-\tau_{2}}\big(y_{(1)}\big) \# v_{\tau_{1}}\big(y_{(2)}\big),
\end{gather*}
where $u_{\gamma^{-1}}\colon T_{m^{2}}(q) \to T_{m^{2}}(q)$, $r_{m-\tau_{2}}\colon K[C_{n}] \to T_{m^{2}}(q)$ and $v_{\tau_{1}}\colon K[C_{n}] \to K[C_{n}]$ are given as follows for all $i, j \in \{0, 1,\dots, m-1\}$ and $k \in \{0, 1,\dots, n-1\}$
\begin{gather*}
u_{\gamma^{-1}}\big(h^{i}x^{j}\big) = \gamma^{-j} h^{i}x^{j}, \qquad r_{m-\tau_{2}}\big(g^{k}\big) = h^{k(m-\tau_{2})}, \qquad v_{\tau_{1}}\big(g^{k}\big) = g^{k \tau_{1}}.
\end{gather*}
The proof will be finished once we show that $\psi_{(u_{\gamma^{-1}}, \varepsilon, r_{m-\tau_{2}}, v_{\tau_{1}})}$ is the
inverse of $\psi_{(u_{\gamma}, \varepsilon, r_{l}, v_{s})}$. Now for
all $i, j \in \{0, 1,\dots, m-1\}$ and $k \in \{0, 1,\dots, n-1\}$ we have
\begin{gather*}
\psi_{(u_{\gamma}, \varepsilon, r_{l}, v_{s})} \circ \psi_{(u_{\gamma^{-1}}, \varepsilon, r_{m-\tau_{2}}, v_{\tau_{1}})}\big(h^{i}x^{j} \# g^{k}\big) = \psi_{(u_{\gamma}, \varepsilon, r_{l}, v_{s})} \big(u_{\gamma^{-1}}\big(h^{i}x^{j}\big)r_{m-\tau_{2}}\big(g^{k}\big) \# v_{\tau_{1}}\big(g^{k}\big)\big) \\
\qquad \ {} = \psi_{(u_{\gamma}, \varepsilon, r_{l}, v_{s})} \big(\gamma^{-j} h^{i}x^{j}h^{k (m-\tau_{2})} \# g^{k \tau_{1}}\big)
 = \gamma^{-j} u_{\gamma}\big(h^{i}x^{j}h^{k (m-\tau_{2})}\big) r_{l}\big(g^{k \tau_{1}}\big) \# v_{s}\big(g^{k \tau_{1}}\big) \\
 \qquad \ {} = \gamma^{-j} \gamma^{j} h^{i}x^{j} h^{k (m-\tau_{2})} h^{k \tau_{1} l} \# g^{k \tau_{1} s}
 = h^{i}x^{j} h^{k (\underline{\tau_{1}} l -\underline{ \tau_{2}})} \# g^{k \underline{\tau_{1}} s} \\
\qquad{} \stackrel{\eqref{impr}} {=} h^{i}x^{j} h^{k(\tau l - \alpha n l - l \tau + \beta m)} \# g^{ks(\tau - \alpha n)}
 = h^{i}x^{j} h^{-k \alpha n l } \# g^{ks\tau} \\
 \qquad \ {} = h^{i}x^{j} h^{-k \alpha n l } \# g^{k(1-n \mu)} \stackrel{m \,|\, nl} = h^{i}x^{j} \# g^{k}.
\end{gather*}
Similarly one can prove that $\psi_{(u_{\gamma^{-1}}, \varepsilon, r_{m-\tau_{2}}, v_{\tau_{1}})} \circ \psi_{(u_{\gamma}, \varepsilon, r_{l}, v_{s})} (h^{i}x^{j} \# g^{k}) = h^{i}x^{j} \# g^{k}$ for all $i, j \in \{0, 1,\dots, m-1\}$ and $k \in \{0, 1,\dots, n-1\}$ and the proof is now finished.
\end{proof}
